\documentclass[11pt,a4paper]{article}

% ----------------------------------------------------------------------
% Packages
% ----------------------------------------------------------------------
\usepackage[T1]{fontenc}
\usepackage[utf8]{inputenc}
\usepackage[english]{babel}
\usepackage{lmodern}
\usepackage{microtype}
\usepackage[margin=2.5cm]{geometry}
\usepackage{amsmath,amssymb,mathtools}
\usepackage{enumitem}
\usepackage{xcolor}
\usepackage{booktabs}
\usepackage{array}
\usepackage[colorlinks=true,linkcolor=blue!60!black,urlcolor=blue!60!black]{hyperref}

\setlength{\parindent}{0pt}
\setlength{\parskip}{0.6em}
\setlist{itemsep=0.2em, topsep=0.3em}

\newcommand{\Veff}{V_{\text{eff}}}

\title{Derivation of the Schwarzschild Equations of Motion\\(Hamiltonian Formalism)}
\author{Nugie Saputra}
\date{\today}

\begin{document}

\maketitle

\begin{abstract}
This document derives the equations of motion for geodesics of test particles and photons in a static, spherically symmetric spacetime using Hamiltonian mechanics. It ends with a system of first-order ordinary differential equations that can be integrated numerically with the fourth-order Runge--Kutta method. The implementation is in the \texttt{schwarzschild} Python package of this repository.
\end{abstract}

\tableofcontents

% ======================================================================
\section{The covariant metric \texorpdfstring{($g_{\mu\nu}$)}{(g\_mu nu)}}
\label{sec:metric}

For a static spacetime with spherical symmetry the line element $ds^2$ is
\begin{equation*}
ds^2 = -f(r)\,dt^2 + \frac{dr^2}{f(r)} + r^2 \left( d\theta^2 + \sin^2\theta\,d\phi^2 \right) \tag{1}\label{eq:1}
\end{equation*}

From equation~\eqref{eq:1} the covariant metric $g_{\mu\nu}$ is diagonal with non-zero components
\begin{align*}
    g_{tt} &= -f(r) \tag{2a}\label{eq:2a} \\
    g_{rr} &= \frac{1}{f(r)} \tag{2b} \\
    g_{\theta\theta} &= r^2 \tag{2c} \\
    g_{\phi\phi} &= r^2 \sin^2\theta \tag{2d}
\end{align*}

% ======================================================================
\section{The contravariant metric \texorpdfstring{($g^{\mu\nu}$)}{(g\^{}mu nu)}}

The contravariant metric $g^{\mu\nu}$ is the inverse of $g_{\mu\nu}$. Since $g_{\mu\nu}$ is diagonal, each contravariant component is simply the reciprocal of the corresponding covariant component in equation~(2):
\begin{align*}
    g^{tt} &= \frac{1}{g_{tt}} = -\frac{1}{f(r)} \tag{3a}\label{eq:3a} \\
    g^{rr} &= \frac{1}{g_{rr}} = f(r) \tag{3b} \\
    g^{\theta\theta} &= \frac{1}{g_{\theta\theta}} = \frac{1}{r^2} \tag{3c} \\
    g^{\phi\phi} &= \frac{1}{g_{\phi\phi}} = \frac{1}{r^2 \sin^2\theta} \tag{3d}
\end{align*}

% ======================================================================
\section{The Hamiltonian}

In General Relativity the Hamiltonian of a particle (or photon) is
\begin{equation*}
H = \frac{1}{2} g^{\mu\nu} p_\mu p_\nu \tag{4}\label{eq:4}
\end{equation*}

Substituting the contravariant components from equation~(3) into equation~\eqref{eq:4} gives
\begin{equation*}
H = \frac{1}{2} \left[ -\frac{p_t^2}{f(r)} + f(r) p_r^2 + \frac{p_\theta^2}{r^2} + \frac{p_\phi^2}{r^2 \sin^2\theta} \right] \tag{5}\label{eq:5}
\end{equation*}

% ======================================================================
\section{Hamilton's equations and conserved quantities}

The momenta conjugate to the coordinates obey the canonical equation, with $\lambda$ an affine parameter:
\begin{equation*}
\dot{p}_\mu = -\frac{\partial H}{\partial x^\mu} \tag{6}\label{eq:6}
\end{equation*}

The Hamiltonian~\eqref{eq:5} does not depend explicitly on the time coordinate $t$ or on the azimuthal angle $\phi$ (cyclic coordinates). Consequently:
\begin{itemize}
  \item \textbf{Conservation of energy ($E$):}
  \begin{equation*}
  \frac{\partial H}{\partial t} = 0 \implies \dot{p}_t = 0 \implies p_t = -E \quad (\text{constant}) \tag{7a}\label{eq:7a}
  \end{equation*}
  This constant of motion is defined as the negative of the specific energy of the particle.

  \item \textbf{Conservation of angular momentum ($L$):}
  \begin{equation*}
  \frac{\partial H}{\partial \phi} = 0 \implies \dot{p}_\phi = 0 \implies p_\phi = L \quad (\text{constant}) \tag{7b}\label{eq:7b}
  \end{equation*}
  This constant of motion is the specific angular momentum about the symmetry axis.
\end{itemize}

% ======================================================================
\section{The equatorial plane}
\label{sec:equatorial}

Because the spacetime is fully spherically symmetric, a particle's trajectory always stays in a single plane. Without loss of generality the coordinates can be oriented so that the motion takes place in the equatorial plane:
\begin{equation*}
\theta = \frac{\pi}{2}, \quad p_\theta = 0, \quad \sin\theta = 1
\end{equation*}

Substituting this condition and the conserved quantities~\eqref{eq:7a} and~\eqref{eq:7b} into equation~\eqref{eq:5}, the Hamiltonian reduces to
\begin{align*}
    H &= \frac{1}{2} \left[ -\frac{(-E)^2}{f(r)} + f(r) p_r^2 + 0 + \frac{L^2}{r^2(1)} \right] \\
    H &= \frac{1}{2} \left[ -\frac{E^2}{f(r)} + f(r) p_r^2 + \frac{L^2}{r^2} \right] \tag{8}\label{eq:8}
\end{align*}

% ======================================================================
\section{Momentum and radial velocity}

The second canonical equation relates the coordinate velocity to the derivative of the Hamiltonian with respect to the conjugate momentum:
\begin{equation*}
\dot{x}^\mu = \frac{\partial H}{\partial p_\mu} \tag{9}\label{eq:9}
\end{equation*}

For the radial component,
\begin{align*}
    \dot{r} &= \frac{\partial H}{\partial p_r} = \frac{1}{2} \left[ 2 f(r) p_r \right] \\
    \dot{r} &= f(r) p_r \tag{10}\label{eq:10}
\end{align*}

The radial momentum can therefore be written in terms of the radial velocity:
\begin{equation*}
p_r = \frac{\dot{r}}{f(r)} \tag{11}\label{eq:11}
\end{equation*}

% ======================================================================
\section{The radial equation of motion}

From the normalization of the 4-momentum ($g^{\mu\nu} p_\mu p_\nu = -\epsilon$), the constant value of the Hamiltonian is
\begin{equation*}
H = -\frac{\epsilon}{2} \tag{12}\label{eq:12}
\end{equation*}
where
\begin{itemize}
  \item $\epsilon = 1$ for massive particles (timelike geodesics),
  \item $\epsilon = 0$ for photons and other massless particles (null geodesics).
\end{itemize}

Substitute equations~\eqref{eq:11} and~\eqref{eq:12} into equation~\eqref{eq:8}:
\begin{align*}
    -\frac{\epsilon}{2} &= \frac{1}{2} \left[ -\frac{E^2}{f(r)} + f(r) \left( \frac{\dot{r}}{f(r)} \right)^2 + \frac{L^2}{r^2} \right] \\
    -\epsilon &= -\frac{E^2}{f(r)} + f(r) \frac{\dot{r}^2}{f(r)^2} + \frac{L^2}{r^2} \\
    -\epsilon &= -\frac{E^2}{f(r)} + \frac{\dot{r}^2}{f(r)} + \frac{L^2}{r^2}
\end{align*}

Multiply both sides by $f(r)$:
\begin{equation*}
-\epsilon f(r) = -E^2 + \dot{r}^2 + f(r) \frac{L^2}{r^2}
\end{equation*}

Move terms to isolate $\dot{r}^2$:
\begin{align*}
    \dot{r}^2 &= E^2 - \epsilon f(r) - f(r) \frac{L^2}{r^2} \\
    \dot{r}^2 &= E^2 - f(r) \left( \epsilon + \frac{L^2}{r^2} \right) \tag{13}\label{eq:13}
\end{align*}

% ======================================================================
\section{The effective potential}

The radial equation~\eqref{eq:13} can be written like the conservation of mechanical energy in classical mechanics:
\begin{equation*}
\dot{r}^2 + \Veff(r) = E^2 \tag{14}\label{eq:14}
\end{equation*}
by defining the \textbf{effective potential} $\Veff(r)$ as
\begin{equation*}
\Veff(r) = f(r) \left( \epsilon + \frac{L^2}{r^2} \right) \tag{15}\label{eq:15}
\end{equation*}

\subsection*{Special case: the Schwarzschild metric}

Outside a non-rotating body of mass $M$ (the Schwarzschild metric) the metric function is
\begin{equation*}
f(r) = 1 - \frac{2M}{r} \tag{16}\label{eq:16}
\end{equation*}
in geometrized units where $G = c = 1$.

Substituting $f(r)$ into equation~\eqref{eq:15} gives the explicit effective potential
\begin{align*}
    \Veff(r) &= \left( 1 - \frac{2M}{r} \right) \left( \epsilon + \frac{L^2}{r^2} \right) \\
    \Veff(r) &= \epsilon - \frac{2M\epsilon}{r} + \frac{L^2}{r^2} - \frac{2ML^2}{r^3} \tag{17}\label{eq:17}
\end{align*}

Physical interpretation of the terms:
\begin{enumerate}
  \item $\epsilon$: the rest energy of the particle ($1$ for massive particles, $0$ for photons).
  \item $-\frac{2M\epsilon}{r}$: the Newtonian gravitational potential.
  \item $\frac{L^2}{r^2}$: the Newtonian centrifugal barrier.
  \item $-\frac{2ML^2}{r^3}$: the General Relativity correction, responsible for the precession of the orbit's periapsis and for the strong attraction at short distances.
\end{enumerate}

% ======================================================================
\section{Why use the full Hamilton equations?}

The radial equation~\eqref{eq:13} has the form $\dot{r}^2 = E^2 - \Veff(r)$. For numerical work this form is \textbf{not ideal}:
\begin{enumerate}
  \item Obtaining $\dot{r}$ requires a square root, $\dot{r} = \pm\sqrt{E^2 - \Veff}$, so the $\pm$ sign has to be switched by hand each time the particle passes a turning point ($\dot{r} = 0$), for example at periapsis or apoapsis.
  \item Near a turning point $E^2 - \Veff \to 0$; a small numerical error can make the argument of the square root negative and the integration fails.
\end{enumerate}

The solution is to use \textbf{all} of Hamilton's canonical equations~\eqref{eq:6} and~\eqref{eq:9}. They form a system of \textbf{first-order} ordinary differential equations (ODEs) that has no square root and passes through turning points automatically:
\begin{equation*}
\boxed{\;\dot{x}^\mu = \frac{\partial H}{\partial p_\mu}, \qquad \dot{p}_\mu = -\frac{\partial H}{\partial x^\mu}\;} \tag{18}\label{eq:18}
\end{equation*}
where a dot denotes the derivative with respect to the affine parameter $\lambda$ (for massive particles $\lambda = \tau$ is the proper time). The 4 coordinates $(t, r, \theta, \phi)$ and 4 momenta $(p_t, p_r, p_\theta, p_\phi)$ give 8 first-order ODEs.

As a starting point, the general Hamiltonian~\eqref{eq:5} is rewritten:
\begin{equation*}
H = \frac{1}{2} \left[ -\frac{p_t^2}{f(r)} + f(r)\, p_r^2 + \frac{p_\theta^2}{r^2} + \frac{p_\phi^2}{r^2 \sin^2\theta} \right]
\end{equation*}

% ======================================================================
\section{Velocity equations \texorpdfstring{($\dot{x}^\mu = \partial H / \partial p_\mu$)}{}}

Each momentum $p_\mu$ appears in exactly one term inside the brackets of $H$, so its partial derivative follows immediately.

\textbf{The $t$ component.} Only the first term depends on $p_t$:
\begin{equation*}
\dot{t} = \frac{\partial H}{\partial p_t} = \frac{1}{2}\left( -\frac{2 p_t}{f(r)} \right) = -\frac{p_t}{f(r)} \tag{19a}\label{eq:19a}
\end{equation*}

\textbf{The $r$ component.} Only the second term depends on $p_r$ (same as equation~\eqref{eq:10}):
\begin{equation*}
\dot{r} = \frac{\partial H}{\partial p_r} = \frac{1}{2}\left( 2 f(r)\, p_r \right) = f(r)\, p_r \tag{19b}\label{eq:19b}
\end{equation*}

\textbf{The $\theta$ component.} Only the third term depends on $p_\theta$:
\begin{equation*}
\dot{\theta} = \frac{\partial H}{\partial p_\theta} = \frac{1}{2}\left( \frac{2 p_\theta}{r^2} \right) = \frac{p_\theta}{r^2} \tag{19c}\label{eq:19c}
\end{equation*}

\textbf{The $\phi$ component.} Only the fourth term depends on $p_\phi$:
\begin{equation*}
\dot{\phi} = \frac{\partial H}{\partial p_\phi} = \frac{1}{2}\left( \frac{2 p_\phi}{r^2 \sin^2\theta} \right) = \frac{p_\phi}{r^2 \sin^2\theta} \tag{19d}\label{eq:19d}
\end{equation*}

Note that equation~(19) is just $\dot{x}^\mu = g^{\mu\nu} p_\nu = p^\mu$, the relation between covariant and contravariant momenta (raising an index with the metric).

% ======================================================================
\section{Momentum equations \texorpdfstring{($\dot{p}_\mu = -\partial H / \partial x^\mu$)}{}}

\subsection{The \texorpdfstring{$t$ and $\phi$}{t and phi} components (cyclic coordinates)}

Because $H$ does not contain $t$ or $\phi$ explicitly,
\begin{align*}
    \dot{p}_t &= -\frac{\partial H}{\partial t} = 0 \tag{20a}\label{eq:20a} \\
    \dot{p}_\phi &= -\frac{\partial H}{\partial \phi} = 0 \tag{20b}\label{eq:20b}
\end{align*}
which recovers the conserved quantities $p_t = -E$ and $p_\phi = L$ of equation~(7).

\subsection{The \texorpdfstring{$r$}{r} component}

The coordinate $r$ appears in all four terms. Write $f'(r) \equiv \dfrac{df}{dr}$ and differentiate term by term:
\begin{itemize}
  \item First term (use $\frac{d}{dr}\left(\frac{1}{f}\right) = -\frac{f'}{f^2}$):
  \begin{equation*}
  \frac{\partial}{\partial r}\left( -\frac{p_t^2}{f} \right) = -p_t^2 \left( -\frac{f'}{f^2} \right) = \frac{f'\, p_t^2}{f^2}
  \end{equation*}
  \item Second term:
  \begin{equation*}
  \frac{\partial}{\partial r}\left( f\, p_r^2 \right) = f'\, p_r^2
  \end{equation*}
  \item Third term:
  \begin{equation*}
  \frac{\partial}{\partial r}\left( \frac{p_\theta^2}{r^2} \right) = -\frac{2 p_\theta^2}{r^3}
  \end{equation*}
  \item Fourth term:
  \begin{equation*}
  \frac{\partial}{\partial r}\left( \frac{p_\phi^2}{r^2 \sin^2\theta} \right) = -\frac{2 p_\phi^2}{r^3 \sin^2\theta}
  \end{equation*}
\end{itemize}

Add them and multiply by $\tfrac{1}{2}$:
\begin{equation*}
\frac{\partial H}{\partial r} = \frac{f'}{2}\left( \frac{p_t^2}{f^2} + p_r^2 \right) - \frac{p_\theta^2}{r^3} - \frac{p_\phi^2}{r^3 \sin^2\theta}
\end{equation*}
Hence
\begin{equation*}
\dot{p}_r = -\frac{\partial H}{\partial r} = -\frac{f'(r)}{2}\left( \frac{p_t^2}{f(r)^2} + p_r^2 \right) + \frac{p_\theta^2}{r^3} + \frac{p_\phi^2}{r^3 \sin^2\theta} \tag{20c}\label{eq:20c}
\end{equation*}

\subsection{The \texorpdfstring{$\theta$}{theta} component}

Only the fourth term depends on $\theta$. With $\frac{d}{d\theta}\left(\sin^{-2}\theta\right) = -2\sin^{-3}\theta \cos\theta$,
\begin{equation*}
\frac{\partial H}{\partial \theta} = \frac{1}{2}\, \frac{p_\phi^2}{r^2} \left( -\frac{2\cos\theta}{\sin^3\theta} \right) = -\frac{p_\phi^2 \cos\theta}{r^2 \sin^3\theta}
\end{equation*}
Hence
\begin{equation*}
\dot{p}_\theta = -\frac{\partial H}{\partial \theta} = \frac{p_\phi^2 \cos\theta}{r^2 \sin^3\theta} \tag{20d}\label{eq:20d}
\end{equation*}

\subsection{Explicit form for Schwarzschild}

For $f(r) = 1 - \dfrac{2M}{r}$ (equation~\eqref{eq:16}) the derivative is
\begin{equation*}
f'(r) = \frac{2M}{r^2} \tag{21}\label{eq:21}
\end{equation*}
and the combination appearing in~\eqref{eq:20c} simplifies:
\begin{equation*}
\frac{f'}{f^2} = \frac{2M/r^2}{\left(1 - \frac{2M}{r}\right)^2} = \frac{2M/r^2}{(r - 2M)^2 / r^2} = \frac{2M}{(r - 2M)^2}
\end{equation*}
Substituting into~\eqref{eq:20c} gives
\begin{equation*}
\dot{p}_r = -\frac{M\, p_t^2}{(r - 2M)^2} - \frac{M\, p_r^2}{r^2} + \frac{p_\theta^2}{r^3} + \frac{p_\phi^2}{r^3 \sin^2\theta} \tag{22}\label{eq:22}
\end{equation*}

\subsection{Summary: the full 8-equation system (3D)}

\begin{equation*}
\begin{aligned}
    \dot{t} &= -\frac{p_t}{1 - 2M/r}, &
    \dot{p}_t &= 0, \\[4pt]
    \dot{r} &= \left(1 - \frac{2M}{r}\right) p_r, &
    \dot{p}_r &= -\frac{M\, p_t^2}{(r - 2M)^2} - \frac{M\, p_r^2}{r^2} + \frac{p_\theta^2}{r^3} + \frac{p_\phi^2}{r^3 \sin^2\theta}, \\[4pt]
    \dot{\theta} &= \frac{p_\theta}{r^2}, &
    \dot{p}_\theta &= \frac{p_\phi^2 \cos\theta}{r^2 \sin^3\theta}, \\[4pt]
    \dot{\phi} &= \frac{p_\phi}{r^2 \sin^2\theta}, &
    \dot{p}_\phi &= 0.
\end{aligned} \tag{23}\label{eq:23}
\end{equation*}

System~\eqref{eq:23} holds for motion in an arbitrary plane and can be used directly to simulate 3D motion (for example many photons with different orientations for ray tracing). It is singular at $\sin\theta = 0$ (the polar axis), however, so for a single particle it is more practical to use the reduction to the equatorial plane below.

% ======================================================================
\section{Reduction to the equatorial plane: the system used in simulations}

\subsection{Consistency of the equatorial plane}

Choose the initial conditions $\theta(0) = \pi/2$ and $p_\theta(0) = 0$. From~\eqref{eq:23},
\begin{equation*}
\dot{\theta}\big|_{\lambda=0} = \frac{0}{r^2} = 0, \qquad
\dot{p}_\theta\big|_{\lambda=0} = \frac{L^2 \cos(\pi/2)}{r^2 \sin^3(\pi/2)} = \frac{L^2 \cdot 0}{r^2} = 0
\end{equation*}

Since neither $\theta$ nor $p_\theta$ changes, they \textbf{remain} $\pi/2$ and $0$ along the whole trajectory. The motion is indeed confined to the equatorial plane, which justifies the assumption of Section~\ref{sec:equatorial}.

\subsection{The reduced system}

Substitute $\theta = \pi/2$, $p_\theta = 0$, $p_t = -E$ and $p_\phi = L$ into~\eqref{eq:23}. The equations for $\theta$, $p_\theta$, $p_t$ and $p_\phi$ become trivial, leaving \textbf{four first-order ODEs}:
\begin{equation*}
\boxed{
\begin{aligned}
    \frac{dt}{d\lambda} &= \frac{E}{1 - \dfrac{2M}{r}} \\[6pt]
    \frac{dr}{d\lambda} &= \left(1 - \frac{2M}{r}\right) p_r \\[6pt]
    \frac{d\phi}{d\lambda} &= \frac{L}{r^2} \\[6pt]
    \frac{dp_r}{d\lambda} &= -\frac{M E^2}{(r - 2M)^2} - \frac{M\, p_r^2}{r^2} + \frac{L^2}{r^3}
\end{aligned}
} \tag{24}\label{eq:24}
\end{equation*}

Equation~\eqref{eq:24} is the \textbf{final result, ready for numerical simulation}. Some important remarks:
\begin{itemize}
  \item \textbf{The parameter $\epsilon$ does not appear explicitly in~\eqref{eq:24}.} Massive particles ($\epsilon = 1$) and photons ($\epsilon = 0$) obey exactly the same equations; the difference enters only through the \textbf{initial condition} $p_r(0)$ (see Section~\ref{sec:initial}).
  \item The right-hand side of~\eqref{eq:24} \textbf{does not depend} on $t$, $\phi$ or $\lambda$ (an autonomous system). Only $r$ and $p_r$ determine the rates of change; $t$ and $\phi$ are simply ``accumulated''.
  \item Physical meaning of the terms in $\dot{p}_r$:
  \begin{itemize}
    \item $-\dfrac{ME^2}{(r-2M)^2}$: the gravitational pull, affected by the energy (it diverges at the horizon),
    \item $-\dfrac{M p_r^2}{r^2}$: a correction from the curvature of the radial part of the metric,
    \item $+\dfrac{L^2}{r^3}$: the centrifugal force.
  \end{itemize}
\end{itemize}

\subsection{General form (arbitrary \texorpdfstring{$f(r)$}{f(r)})}

To change the metric (for example Reissner--Nordstr\"om, de~Sitter--Schwarzschild, etc.), use the general form before substituting $f$:
\begin{equation*}
\begin{aligned}
    \dot{t} &= \frac{E}{f(r)}, &
    \dot{r} &= f(r)\, p_r, \\[4pt]
    \dot{\phi} &= \frac{L}{r^2}, &
    \dot{p}_r &= -\frac{f'(r)}{2}\left( \frac{E^2}{f(r)^2} + p_r^2 \right) + \frac{L^2}{r^3}
\end{aligned} \tag{25}\label{eq:25}
\end{equation*}
so the code only needs different functions \texttt{f(r)} and \texttt{df(r)}.

% ======================================================================
\section{Verification: consistency with the effective potential}

To check~\eqref{eq:24}, derive the radial acceleration $\ddot{r}$ from Hamilton's equations and compare it with the result from equation~\eqref{eq:14}.

\textbf{From equation~\eqref{eq:14}.} Differentiate $\dot{r}^2 + \Veff = E^2$ with respect to $\lambda$:
\begin{equation*}
2\dot{r}\ddot{r} + \frac{d\Veff}{dr}\dot{r} = 0 \implies \ddot{r} = -\frac{1}{2}\frac{d\Veff}{dr}
\end{equation*}
With $\Veff = f\left(\epsilon + \frac{L^2}{r^2}\right)$,
\begin{equation*}
\ddot{r} = -\frac{f'}{2}\left(\epsilon + \frac{L^2}{r^2}\right) + \frac{f L^2}{r^3} \tag{26}\label{eq:26}
\end{equation*}

\textbf{From Hamilton's equations.} Differentiate~\eqref{eq:19b} with the chain rule, $\frac{d f}{d\lambda} = f'\dot{r} = f' f p_r$:
\begin{equation*}
\ddot{r} = \frac{d}{d\lambda}\left( f p_r \right) = f' \dot{r}\, p_r + f \dot{p}_r = f f' p_r^2 + f \left[ -\frac{f'}{2}\left( \frac{E^2}{f^2} + p_r^2 \right) + \frac{L^2}{r^3} \right]
\end{equation*}
\begin{equation*}
\ddot{r} = \frac{f'}{2}\left( f p_r^2 \right) - \frac{f'}{2}\frac{E^2}{f} + \frac{f L^2}{r^3}
\end{equation*}

The Hamiltonian constraint $H = -\epsilon/2$ (equations~\eqref{eq:8} and~\eqref{eq:12}) gives $f p_r^2 = \dfrac{E^2}{f} - \epsilon - \dfrac{L^2}{r^2}$. Substituting,
\begin{equation*}
\ddot{r} = \frac{f'}{2}\left( \frac{E^2}{f} - \epsilon - \frac{L^2}{r^2} \right) - \frac{f'}{2}\frac{E^2}{f} + \frac{f L^2}{r^3}
= -\frac{f'}{2}\left(\epsilon + \frac{L^2}{r^2}\right) + \frac{f L^2}{r^3}
\end{equation*}

This result is \textbf{identical} to~\eqref{eq:26}. $\checkmark$

For Schwarzschild ($f' = 2M/r^2$),
\begin{equation*}
\ddot{r} = -\frac{M\epsilon}{r^2} + \frac{L^2}{r^3} - \frac{3ML^2}{r^4} \tag{27}\label{eq:27}
\end{equation*}
which is the Newtonian ``force'' plus the relativistic correction $-3ML^2/r^4$.

% ======================================================================
\section{Initial conditions, control quantities and stopping criteria}
\label{sec:initial}

\subsection{Units}

Use geometrized units $G = c = 1$ and choose $M = 1$. Then $r$, $t$, $\lambda$ and $L$ are expressed in units of $M$, while $E$ is dimensionless. The event horizon is at $r_s = 2M = 2$.

\subsection{Determining \texorpdfstring{$p_r(0)$}{p\_r(0)} from the Hamiltonian constraint}

The initial state vector is $\left(t_0, r_0, \phi_0, p_{r,0}\right)$, usually with $t_0 = 0$ and $\phi_0 = 0$. The constants $E$ and $L$ can be chosen freely, but $p_{r,0}$ is \textbf{not} free: it must satisfy $H = -\epsilon/2$. From~\eqref{eq:13} and~\eqref{eq:11},
\begin{equation*}
p_{r,0} = \sigma\, \frac{\sqrt{E^2 - \Veff(r_0)}}{f(r_0)}, \qquad \sigma = \begin{cases} -1 & \text{ingoing} \\ +1 & \text{outgoing} \end{cases} \tag{28}\label{eq:28}
\end{equation*}
provided $E^2 \geq \Veff(r_0)$. The square root is used \textbf{only once}, here; during the integration the sign of $p_r$ is handled automatically by~\eqref{eq:24}.

\subsection{Standard test cases}

\textbf{(a) Circular orbit of a massive particle ($\epsilon = 1$).} The conditions are $\Veff'(r_c) = 0$ and $E^2 = \Veff(r_c)$. From $\Veff' = \frac{2M}{r^2}\left(1 + \frac{L^2}{r^2}\right) - \frac{2 f L^2}{r^3} = 0$, multiplying by $r^4/2$:
\begin{equation*}
M r^2 + M L^2 - L^2 r + 2ML^2 = 0 \implies L^2 (r - 3M) = M r^2
\end{equation*}
\begin{equation*}
L = \sqrt{\frac{M r_c^2}{r_c - 3M}}, \qquad
E = \frac{1 - 2M/r_c}{\sqrt{1 - 3M/r_c}}, \qquad
p_{r,0} = 0 \tag{29}\label{eq:29}
\end{equation*}

Circular orbits are stable only for $r_c \geq 6M$ (the \emph{Innermost Stable Circular Orbit}, ISCO). The numerical result must give a constant $r(\lambda) = r_c$ and the angular frequency $d\phi/dt = \sqrt{M/r_c^3}$ (Kepler's law, which is exact in Schwarzschild coordinates).

\textbf{(b) Bound orbit with precession ($\epsilon = 1$).} Start at apoapsis: choose $r_0$ and $L$, set $p_{r,0} = 0$ and $E = \sqrt{\Veff(r_0)}$. For example $r_0 = 20M$, $L = 4.2M$ gives a ``rosette'' orbit with $r \in [10.33M,\; 20M]$.

\textbf{(c) Photons ($\epsilon = 0$).} Photon paths depend only on the impact parameter $b = L/E$, so $E = 1$ and $L = b$ can be used. The peak of the potential $\Veff = fL^2/r^2$ is at $r = 3M$ (the \emph{photon sphere}) with value $L^2/(27M^2)$, so the critical impact parameter is
\begin{equation*}
b_c = 3\sqrt{3}\, M \approx 5.196\, M \tag{30}\label{eq:30}
\end{equation*}

Photons with $b < b_c$ are captured by the black hole; those with $b > b_c$ are deflected and escape. For a photon coming from far away ($r_0 \gg M$):
\begin{equation*}
p_{r,0} = -\frac{\sqrt{1 - f(r_0)\, b^2 / r_0^2}}{f(r_0)}.
\end{equation*}

\subsection{The Hamiltonian constraint as an error check}

Because $H$ is conserved exactly, the quantity
\begin{equation*}
\mathcal{C}(\lambda) = H + \frac{\epsilon}{2} = \frac{1}{2}\left[ -\frac{E^2}{f(r)} + f(r)\, p_r^2 + \frac{L^2}{r^2} \right] + \frac{\epsilon}{2} \tag{31}\label{eq:31}
\end{equation*}
should be zero along the whole trajectory. In a simulation $|\mathcal{C}|$ is computed at every step as an accuracy indicator. Close to the horizon the term $E^2/f$ grows, so it is better to monitor the relative error $|\mathcal{C}| \,/\, \left(E^2 / 2f\right)$.

\subsection{Stopping criteria}
\label{sec:stop}

The integration stops when any of the following holds:
\begin{itemize}
  \item $r \leq 2M(1 + \delta)$, with $\delta \sim 10^{-3}$: the particle falls into the horizon. Schwarzschild coordinates are singular at $r = 2M$ ($f \to 0$, so $\dot{t} \to \infty$ and $p_r \to \infty$), so the integration cannot continue through the horizon.
  \item $r \geq r_{\text{max}}$: the particle or photon has escaped to infinity.
  \item $\lambda \geq \lambda_{\text{max}}$ or the maximum number of steps is reached.
\end{itemize}

\subsection{Conversion to Cartesian coordinates (visualization)}

In the equatorial plane,
\begin{equation*}
x = r\cos\phi, \qquad y = r\sin\phi \tag{32}\label{eq:32}
\end{equation*}

% ======================================================================
\section{The fourth-order Runge--Kutta method (RK4)}

\subsection{The initial value problem}

RK4 solves a system of first-order ODEs of the form
\begin{equation*}
\frac{d\mathbf{y}}{d\lambda} = \mathbf{F}(\lambda, \mathbf{y}), \qquad \mathbf{y}(\lambda_0) = \mathbf{y}_0 \tag{33}\label{eq:33}
\end{equation*}
where $\mathbf{y}$ is the state vector and $\mathbf{F}$ the right-hand side. The parameter $\lambda$ is discretized with step $h$: $\lambda_n = \lambda_0 + n h$ and $\mathbf{y}_n \approx \mathbf{y}(\lambda_n)$.

\subsection{The basic idea}

Exactly,
\begin{equation*}
\mathbf{y}_{n+1} = \mathbf{y}_n + \int_{\lambda_n}^{\lambda_n + h} \mathbf{F}(\lambda, \mathbf{y}(\lambda))\, d\lambda .
\end{equation*}
Euler's method approximates this integral using only the slope at the start of the interval, $h\,\mathbf{F}(\lambda_n, \mathbf{y}_n)$, so its error is $\mathcal{O}(h)$. RK4 improves on this by evaluating the slope at \textbf{four points} in the interval (start, midpoint twice, end) and taking a weighted average, analogous to Simpson's rule
\begin{equation*}
\int_{\lambda_n}^{\lambda_n + h} F\, d\lambda \approx \frac{h}{6}\left[F_{\text{start}} + 4F_{\text{mid}} + F_{\text{end}}\right],
\end{equation*}
where the midpoint weight $4$ is split into $2 + 2$.

\subsection{The RK4 formulas}

\begin{align*}
    \mathbf{k}_1 &= \mathbf{F}\left(\lambda_n,\; \mathbf{y}_n\right) \tag{34a}\label{eq:34a} \\
    \mathbf{k}_2 &= \mathbf{F}\left(\lambda_n + \tfrac{h}{2},\; \mathbf{y}_n + \tfrac{h}{2}\mathbf{k}_1\right) \tag{34b} \\
    \mathbf{k}_3 &= \mathbf{F}\left(\lambda_n + \tfrac{h}{2},\; \mathbf{y}_n + \tfrac{h}{2}\mathbf{k}_2\right) \tag{34c} \\
    \mathbf{k}_4 &= \mathbf{F}\left(\lambda_n + h,\; \mathbf{y}_n + h\,\mathbf{k}_3\right) \tag{34d} \\[4pt]
    \mathbf{y}_{n+1} &= \mathbf{y}_n + \frac{h}{6}\left( \mathbf{k}_1 + 2\mathbf{k}_2 + 2\mathbf{k}_3 + \mathbf{k}_4 \right) \tag{34e}\label{eq:34e}
\end{align*}

Meaning of each stage:
\begin{itemize}
  \item $\mathbf{k}_1$: the slope at the start of the interval (same as Euler),
  \item $\mathbf{k}_2$: the slope at the midpoint, using a half-step Euler prediction with $\mathbf{k}_1$,
  \item $\mathbf{k}_3$: the slope at the midpoint again, using the improved prediction with $\mathbf{k}_2$,
  \item $\mathbf{k}_4$: the slope at the end of the interval, using a full-step prediction with $\mathbf{k}_3$.
\end{itemize}

\subsection{Error}

The local (per-step) error of RK4 is $\mathcal{O}(h^5)$ and the global error (accumulated over $N \propto 1/h$ steps) is $\mathcal{O}(h^4)$. Halving $h$ therefore reduces the global error by a factor of about $2^4 = 16$. This gives a simple convergence test: run the simulation with $h$ and $h/2$ and compare.

Note: RK4 is \textbf{not} a symplectic integrator, so $H$ drifts slowly over very long integrations (thousands of orbits). For simulations of tens to hundreds of orbits with a reasonable $h$ the drift is very small and can be monitored through $\mathcal{C}(\lambda)$ in~\eqref{eq:31}.

% ======================================================================
\section{Applying RK4 to the Schwarzschild equations of motion}

\subsection{State vector and right-hand side}

From system~\eqref{eq:24}, define
\begin{equation*}
\mathbf{y} = \begin{pmatrix} t \\ r \\ \phi \\ p_r \end{pmatrix}, \qquad
\mathbf{F}(\mathbf{y}) = \begin{pmatrix}
    F^t \\ F^r \\ F^\phi \\ F^{p}
\end{pmatrix} = \begin{pmatrix}
    \dfrac{E}{1 - 2M/r} \\[8pt]
    \left(1 - \dfrac{2M}{r}\right) p_r \\[8pt]
    \dfrac{L}{r^2} \\[8pt]
    -\dfrac{ME^2}{(r-2M)^2} - \dfrac{M p_r^2}{r^2} + \dfrac{L^2}{r^3}
\end{pmatrix} \tag{35}\label{eq:35}
\end{equation*}
with $E$, $L$ and $M$ constant parameters. Since $\mathbf{F}$ does not depend on $\lambda$, the argument $\lambda$ in~(34) can be ignored.

\subsection{One RK4 step in explicit form}
\label{sec:rk4explicit}

Let the state at step $n$ be $(t_n, r_n, \phi_n, p_{r,n})$. Since $\mathbf{F}$ depends only on $r$ and $p_r$, the only intermediate values that need to be computed are $r$ and $p_r$. Define
\begin{equation*}
f(r) = 1 - \frac{2M}{r}, \qquad G(r, p_r) \equiv -\frac{ME^2}{(r-2M)^2} - \frac{M p_r^2}{r^2} + \frac{L^2}{r^3}.
\end{equation*}

\textbf{Stage 1} (at $r_n,\, p_{r,n}$):
\begin{equation*}
\begin{aligned}
    k_1^t &= \frac{E}{f(r_n)}, &
    k_1^r &= f(r_n)\, p_{r,n}, \\[4pt]
    k_1^\phi &= \frac{L}{r_n^2}, &
    k_1^p &= G(r_n, p_{r,n})
\end{aligned}
\end{equation*}

\textbf{Stage 2} (at the midpoint using $\mathbf{k}_1$):
\begin{equation*}
r^{(2)} = r_n + \tfrac{h}{2}k_1^r, \qquad p^{(2)} = p_{r,n} + \tfrac{h}{2}k_1^p
\end{equation*}
\begin{equation*}
\begin{aligned}
    k_2^t &= \frac{E}{f(r^{(2)})}, &
    k_2^r &= f(r^{(2)})\, p^{(2)}, \\[4pt]
    k_2^\phi &= \frac{L}{\left(r^{(2)}\right)^2}, &
    k_2^p &= G\left(r^{(2)}, p^{(2)}\right)
\end{aligned}
\end{equation*}

\textbf{Stage 3} (at the midpoint using $\mathbf{k}_2$):
\begin{equation*}
r^{(3)} = r_n + \tfrac{h}{2}k_2^r, \qquad p^{(3)} = p_{r,n} + \tfrac{h}{2}k_2^p
\end{equation*}
\begin{equation*}
\begin{aligned}
    k_3^t &= \frac{E}{f(r^{(3)})}, &
    k_3^r &= f(r^{(3)})\, p^{(3)}, \\[4pt]
    k_3^\phi &= \frac{L}{\left(r^{(3)}\right)^2}, &
    k_3^p &= G\left(r^{(3)}, p^{(3)}\right)
\end{aligned}
\end{equation*}

\textbf{Stage 4} (at the end of the interval using $\mathbf{k}_3$):
\begin{equation*}
r^{(4)} = r_n + h\,k_3^r, \qquad p^{(4)} = p_{r,n} + h\,k_3^p
\end{equation*}
\begin{equation*}
\begin{aligned}
    k_4^t &= \frac{E}{f(r^{(4)})}, &
    k_4^r &= f(r^{(4)})\, p^{(4)}, \\[4pt]
    k_4^\phi &= \frac{L}{\left(r^{(4)}\right)^2}, &
    k_4^p &= G\left(r^{(4)}, p^{(4)}\right)
\end{aligned}
\end{equation*}

\textbf{Update:}
\begin{align*}
    t_{n+1} &= t_n + \frac{h}{6}\left( k_1^t + 2k_2^t + 2k_3^t + k_4^t \right) \\
    r_{n+1} &= r_n + \frac{h}{6}\left( k_1^r + 2k_2^r + 2k_3^r + k_4^r \right) \\
    \phi_{n+1} &= \phi_n + \frac{h}{6}\left( k_1^\phi + 2k_2^\phi + 2k_3^\phi + k_4^\phi \right) \\
    p_{r,n+1} &= p_{r,n} + \frac{h}{6}\left( k_1^p + 2k_2^p + 2k_3^p + k_4^p \right) \tag{36}\label{eq:36}
\end{align*}

\subsection{Choosing the step size \texorpdfstring{$h$}{h}}

\begin{itemize}
  \item Far from the black hole ($r \gtrsim 6M$) a fixed step $h \sim 0.01$--$0.1\,M$ is already very accurate.
  \item Close to the horizon $p_r \sim 1/(r - 2M)$ and $\dot{p}_r \sim 1/(r - 2M)^2$ grow sharply. With a fixed $h$ the RK4 stages can ``jump'' to $r < 2M$ and produce meaningless values. A simple remedy is to shrink the step in proportion to the distance from the horizon:
  \begin{equation*}
  h_n = \min\left( h_0,\; \alpha\,(r_n - 2M) \right), \qquad \alpha \sim 0.02 \tag{37}\label{eq:37}
  \end{equation*}
  Since $\mathbf{F}$ does not depend on $\lambda$, changing $h$ from step to step does not break the RK4 scheme (just accumulate $\lambda_{n+1} = \lambda_n + h_n$).
\end{itemize}

\subsection{Algorithm}

\begin{enumerate}
  \item Set $M$, $\epsilon$, $E$, $L$, $r_0$, $\phi_0$, $t_0$, $h_0$, $\alpha$, $\delta$, $r_{\text{max}}$ and the maximum number of steps $N$.
  \item Compute $p_{r,0}$ from~\eqref{eq:28} (or from~\eqref{eq:29}/\eqref{eq:30} for the test cases).
  \item Form $\mathbf{y}_0 = (t_0, r_0, \phi_0, p_{r,0})$.
  \item For $n = 0, 1, \ldots, N-1$:
  \begin{enumerate}
    \item Compute $h_n$ from~\eqref{eq:37}.
    \item Compute $\mathbf{k}_1, \mathbf{k}_2, \mathbf{k}_3, \mathbf{k}_4$ from~\eqref{eq:35} as in Section~\ref{sec:rk4explicit}.
    \item Update $\mathbf{y}_{n+1}$ with~\eqref{eq:36} and $\lambda_{n+1} = \lambda_n + h_n$.
    \item Store $\mathbf{y}_{n+1}$ and compute $\mathcal{C}$ from~\eqref{eq:31}.
    \item Stop if a stopping criterion (Section~\ref{sec:stop}) is met.
  \end{enumerate}
  \item Convert $(r, \phi)$ to $(x, y)$ with~\eqref{eq:32} for visualization.
\end{enumerate}

\subsection{Implementation}

The algorithm is implemented in the \texttt{schwarzschild} package (directory \texttt{src/schwarzschild/}); the table maps the equations of this document to the code.

\begin{center}
\begin{tabular}{@{}l>{\raggedright\arraybackslash}p{9.3cm}@{}}
\toprule
Equations & Code \\
\midrule
(15) & \texttt{metric.effective\_potential} \\
(16), (21) & \texttt{metric.Schwarzschild.f}, \texttt{.df} \\
(25) (equals (24) for Schwarzschild) & \texttt{equations.rhs} \\
(28) & \texttt{initial\_conditions.initial\_radial\_momentum} \\
(29) & \texttt{analytic.circular\_orbit\_constants} \\
(30) & \texttt{analytic.critical\_impact\_parameter} \\
(31) & \texttt{equations.hamiltonian\_constraint}, \texttt{equations.relative\_constraint\_error} \\
(32) & \texttt{integrators.Trajectory.xy} \\
(34), (36) & \texttt{integrators.rk4\_step} \\
(37), Section~\ref{sec:stop} & \texttt{integrators.integrate} \\
\bottomrule
\end{tabular}
\end{center}

\subsection{Validation}

The test suite (\texttt{tests/}) and the notebook compare the numerical results with closed-form predictions:
\begin{itemize}
  \item The circular orbit $r_c = 10M$ stays at $r = 10M$, and the numerical $d\phi/dt = 0.0316227766$ equals $\sqrt{M/r_c^3}$.
  \item The bound orbit oscillates in $r \in [10.3308M,\; 20M]$ with $|\mathcal{C}| \sim 10^{-16}$. Its periapsis advance per orbit, $2.127094008$~rad, agrees with the exact elliptic-integral result $2.127093985$~rad (below) to a relative difference of $10^{-8}$.
  \item The photon with $b = 5.19M$ is captured, while $b = 5.20M$ escapes after circling the photon sphere down to $r_{\min} \approx 3.07M$, consistent with $b_c = 3\sqrt{3}M \approx 5.196M$ in~\eqref{eq:30}.
  \item For large $b$ the deflection angle follows the weak-field series $\alpha \approx 4M/b + \frac{15\pi}{4}(M/b)^2 + \frac{128}{3}(M/b)^3 + \frac{3465\pi}{64}(M/b)^4$; for $b = 40M$ the numerical value is $0.108104$ against $0.108096$ from the series.
  \item The measured global convergence order of RK4 is $3.996$ (theory: $4$).
\end{itemize}

\textbf{Exact periapsis advance.} With $u = 1/r$, the orbit equation of a massive particle is $(du/d\phi)^2 = 2M(u - u_1)(u - u_2)(u - u_3)$. The turning points $u_1 = 1/r_{\text{apo}}$ and $u_2 = 1/r_{\text{peri}}$ fix the third root through $u_1 + u_2 + u_3 = 1/(2M)$, and the angle swept in one radial period is
\begin{equation*}
\Delta\phi = \frac{4\,K(k^2)}{\sqrt{2M(u_3 - u_1)}}, \qquad k^2 = \frac{u_2 - u_1}{u_3 - u_1},
\end{equation*}
where $K$ is the complete elliptic integral of the first kind. The periapsis advance per orbit is $\Delta\phi - 2\pi$ (function \texttt{analytic.precession\_per\_orbit}).

\end{document}
